\section{More on Matrices and Linear Equations}

Let $A \in \M_{m \times n}(K)$, and consider the homogeneous system of linear equations $A \cdot x = 0$. We denote the space of solutions of this system by:
\begin{equation*}
    \operatorname{Sol}(A, 0) := \{x \in K^n \mid A x = 0\} = \ker(T_A)
\end{equation*}
Note that by the Rank-Nullity Theorem, we have 
\[\dim \operatorname{Sol}(A, 0) = n - \rank(A),\] 
because $\dim \ker(T_A) = n - \dim \operatorname{Im}(T_A)$.

Let $A'$ be a row-reduced echelon matrix that is row-equivalent to $A$ (for example, $A'$ is obtained from $A$ by Gauss elimination). We then know that the number of pivots of $A'$ equals $\rank(A)$, and the number of free variables is $n - \rank(A)$.

We can visualize the structure of $A'$ and its corresponding variables $x_1, \dots, x_n$ as follows. The columns containing the leading $1$'s (pivots) correspond to the pivot variables, while the remaining columns correspond to the free variables. Clearly, the solution spaces are identical:
\begin{equation*}
    \operatorname{Sol}(A, 0) = \operatorname{Sol}(A', 0)
\end{equation*}

Denote by $x_{i_1}, \dots, x_{i_\ell}$ the free variables, where $1 \le i_1 < \dots < i_\ell \le n$ and $\ell = n - \rank(A)$. Similarly, denote by $x_{j_1}, \dots, x_{j_r}$ the pivot variables, where $1 \le j_1 < \dots < j_r \le n$ and $r = \rank(A)$.

For each choice of values for the free variables $x_{i_1}, \dots, x_{i_\ell} \in K$, the pivot variables $x_{j_1}, \dots, x_{j_r}$ are uniquely determined. So, we obtain a map $\Phi: K^\ell \to \operatorname{Sol}(A, 0)$, where $\ell = n - \rank(A)$, defined by:
\begin{equation*}
    \Phi(\lambda_1, \dots, \lambda_\ell) = \text{the unique solution of } Ax = 0 \text{ (which is the solution of } A'x = 0 \text{)}
\end{equation*}
with $x_{i_1} = \lambda_1, \dots, x_{i_k} = \lambda_k, \dots, x_{i_\ell} = \lambda_\ell$.

\begin{proposition}[Linearity and Isomorphism of the Solution Map]
% prop:17.d.1
\label{prop:solution_map_isomorphism}
The map $\Phi$ is linear. Moreover, it is an isomorphism.
\end{proposition}

\begin{proof}
\textbf{Linearity of $\Phi$:} The $k$-th non-zero row of $A'$ gives the following equation:
\begin{equation*}
    x_{j_k} + (\text{linear combination of the free variables that come after } j_k) = 0
\end{equation*}
This allows us to isolate the pivot variable $x_{j_k}$:
\begin{equation*}
    x_{j_k} = - \sum_{\{q \mid 1 \le q \le \ell, \, i_q > j_k\}} a'_{k, i_q} \cdot x_{i_q}
\end{equation*}
This explicitly shows that the pivot variables $x_{j_1}, \dots, x_{j_r}$ depend linearly on the free variables $x_{i_1}, \dots, x_{i_\ell}$. This proves the linearity of $\Phi$.

\textbf{Isomorphism:} To prove that $\Phi$ is an isomorphism, it is enough to show that $\ker(\Phi) = \{0\}$, because $\Phi: K^\ell \to \operatorname{Sol}(A, 0)$ is a linear map between two spaces of the exact same dimension $\ell$.

But this clearly holds: if $\lambda = (\lambda_1, \dots, \lambda_\ell) \in \ker(\Phi)$, which means $\Phi(\lambda_1, \dots, \lambda_\ell) = 0$, then by the definition of $\Phi$, we must have $\lambda_1 = \dots = \lambda_\ell = 0$. This is simply because each $\lambda_k$ is an actual entry in the output vector $\Phi(\lambda_1, \dots, \lambda_\ell)$.

Furthermore, we can explicitly state the inverse map $\Phi^{-1}$. For any $x \in \operatorname{Sol}(A, 0)$, we extract the entries corresponding to the free variables:
\begin{equation*}
    \Phi^{-1}(x) = \begin{pmatrix} x_{i_1} \\ \vdots \\ x_{i_\ell} \end{pmatrix}
\end{equation*}
This completes the proof.
\end{proof}

\begin{corollary}[Structure of Non-Homogeneous Solutions]
% cor:17.d.2
\label{cor:non_homogeneous_solutions}
Let $b \in K^m$, and $A \in \M_{m \times n}(K)$. Then:
\begin{enumerate}
    \item $\operatorname{Sol}(A, b) := \{x \in K^n \mid A \cdot x = b\} \neq \emptyset$ if and only if $b \in \operatorname{ColS}(A)$.
    \item If $\operatorname{Sol}(A, b) \neq \emptyset$ and $y \in \operatorname{Sol}(A, b)$, then:
    \begin{equation*}
        \operatorname{Sol}(A, b) = y + \operatorname{Sol}(A, 0) := \{y + x \mid x \in \operatorname{Sol}(A, 0)\}
    \end{equation*}
\end{enumerate}
\end{corollary}

\begin{proof}
This is left as an exercise.

\textbf{Hints:}
1) Write the matrix $A$ in terms of its column vectors $v_1, \dots, v_n$:
\begin{equation*}
    A = \begin{pmatrix} | & & | \\ v_1 & \dots & v_n \\ | & & | \end{pmatrix}
\end{equation*}
For any vector $x = \begin{pmatrix} x_1 \\ \vdots \\ x_n \end{pmatrix} \in K^n$, the matrix-vector product is a linear combination of the columns:
\begin{equation*}
    A \cdot x = x_1 \begin{pmatrix} | \\ v_1 \\ | \end{pmatrix} + \dots + x_n \begin{pmatrix} | \\ v_n \\ | \end{pmatrix}
\end{equation*}

2) Let $y \in \operatorname{Sol}(A, b)$. First, assume $x \in \operatorname{Sol}(A, 0)$. Then $A \cdot x = 0$ and $A \cdot y = b$. Adding these together yields:
\begin{equation*}
    A(y + x) = A \cdot y + A \cdot x = b + 0 = b
\end{equation*}
Conversely, if $z \in \operatorname{Sol}(A, b)$, meaning $A \cdot z = b$, then define $x := z - y$. We obviously have $z = y + x$, and checking the multiplication gives:
\begin{equation*}
    A \cdot x = A \cdot (z - y) = A \cdot z - A \cdot y = b - b = 0
\end{equation*}
This shows that any solution $z$ can be written as $y + x$ for some $x \in \operatorname{Sol}(A, 0)$.
\end{proof}

\begin{proposition}[Rank Criterion for Solvability]
% prop:17.d.3
\label{prop:rank_solvability_criterion}
Let $A \in \M_{m \times n}(K)$ and $b \in K^m$. The following statements are equivalent:
\begin{enumerate}
    \item $\rank(A|b) = \rank(A)$.
    \item The system of equations $Ax = b$ has a solution.
\end{enumerate}
\end{proposition}

\begin{proof}
Write $A$ in terms of its columns, $A = \begin{pmatrix} | & & | \\ v_1 & \dots & v_n \\ | & & | \end{pmatrix}$. We have the following chain of equivalences:
\begin{align*}
    \rank(A|b) = \rank(A) &\iff \operatorname{col-rank}(A|b) = \operatorname{col-rank}(A) \\
    &\iff \dim \Sp(v_1, \dots, v_n, b) = \dim \Sp(v_1, \dots, v_n) \\
    &\iff b \in \Sp(v_1, \dots, v_n) \\
    &\iff b \in \operatorname{ColS}(A) \\
    &\iff \operatorname{Sol}(A, b) \neq \emptyset
\end{align*}
The last step holds directly by Corollary \ref{cor:non_homogeneous_solutions}, thus finishing the proof.
\end{proof}

\begin{proposition}[Criterion for a Unique Solution]
% prop:17.d.4
\label{prop:unique_solution_rank_criterion}
Let $A \in \M_{m \times n}(K)$ and $b \in K^m$. The following statements are equivalent:
\begin{enumerate}
    \item $\rank(A) = \rank(A|b) = n$.
    \item The system $Ax = b$ has a unique solution.
\end{enumerate}
\end{proposition}

\begin{proof}
This is left as an exercise.

\textbf{Hints:} Write $A = \begin{pmatrix} | & & | \\ v_1 & \dots & v_n \\ | & & | \end{pmatrix}$. Note that:
\begin{equation*}
    \rank(A) = n \iff \dim \Sp(v_1, \dots, v_n) = n
\end{equation*}
Since there are $n$ column vectors, this means that $v_1, \dots, v_n$ are linearly independent, and therefore, they form a basis for $\operatorname{ColS}(A)$.
\end{proof}